\begin{abstract}
Time-lapse microscopy reconstructs cell lineages by linking cells frame to frame, which requires dense imaging that is phototoxic, storage-hungry and often impossible.
We ask how much lineage survives sparse imaging: given two snapshots of a growing bacterial colony taken $H$ minutes apart, assign every cell in the later one to its ancestor in the earlier one.
Between snapshots cells divide several times and colonies deform, so trackers, optimal transport and point-set registration degrade sharply.
We propose \ours, which registers descendant \emph{families} rather than cells: it iteratively moves one predicted family centre per ancestor and scores candidates relative to it.
On One-in-a-Million movies its window-averaged accuracy is 0.96 at a 33-minute gap and 0.47 at 132 minutes, where an ancestor has up to eight descendants; composing its links over 13 of 133 frames recovers two hours of lineage for 0.99 of cells, where off-the-shelf Trackastra, even given true masks, needs 34 frames.
Two routes avoid lineage annotation.
Trained only on a colony simulator, \ours beats every positions-only baseline, optimal transport and forward simulation included, zero-shot on all 27 sequences of four bacterial datasets and, in a pre-registered blind test, on all 28 colonies of three unseen subsets; a plain learned matcher trained the same way does as well, so this transfer is the simulator's.
Chaining one-frame links under a binary-fission prior through unannotated time-lapses yields pseudo-lineage that trains \ours to within 0.02--0.05 of annotation.
Decoded with a single-image genealogy network, this annotation-free pipeline beats a transformer tracker fine-tuned on annotated lineage and decoded the same way by 0.02--0.03 at 33 and 66 minutes (0.07--0.11 against its own output) and matches it at 132.
We also report pre-registered failures: fields of many tiny colonies and steady-state chambers.
\end{abstract}
